\subsection{1-Lipschitz gradient norm preserving network}
In order to build a deep learning classifier based on the hinge-KR optimization problem (Eq.~\eqref{eq:reg_OT}), we have to constrain the Lipschitz constant of the neural network to be equal to 1.  It is known that evaluating it exactly is a NP-hard problem~\cite{NIPS2018_7640}. The simplest way to constraint a network to be 1-Lipschitz is to impose this 1-Lipschitz property to each layer. For dense layers, the initial version of WGAN \cite{Arjovsky2017} consisted of clipping the weights of the layers. However, this is a very crude way to upper-bound Lipschitz constant. Normalizing by the Frobenius norm has also been proposed in~\cite{SalimansK16}. In this paper, we use spectral normalization as proposed in~\cite{Miyato2018SpectralNF}, since the spectral norm is equals to the Lipschitz constant of the layer. At the inference step, we normalize the weights of each layer by the spectral norm of the matrix. This spectral norm is computed by iteratively evaluating the largest singular value with the power iteration algorithm~\cite{GOLUB200035}. This is done during the forward step and taken into account for the gradient computation. In the case of 2D-convolutional layers, normalizing by the spectral norm of convolution kernels is not enough and a supplementary multiplicative constant $\Lambda$ is required (the regularization is then done by dividing $W$ by $\Lambda||W||$). We propose, for zero padding, a tighter estimation of $\Lambda$ than the one proposed in~\cite{cisse_parseval_2017}, computing the average duplication factor of non zero padded values in the feature map:

\begin{equation}
\Lambda=\sqrt{\frac{(k.w-\bar{k}.(\bar{k}+1)).(k.h-\bar{k}.(\bar{k}+1))}{h.w}}
\label{eq:ConvApproxLip}
\end{equation}
for a kernel size equals to $k=2*\bar{k}+1$.
Even if this constant doesn't provide a strict upper bound of the Lipschitz constant (for instance, when the higher values are located in the center of the picture), it behaves very well empirically.
Convolution with stride, pooling layers, detailed explanations and demonstrations are discussed in Appendix~\ref{sec:convStride}.\\

As shown in Property~\ref{coro_norm1}, the optimal function $f^*$ with respect to Eq.~\eqref{eq:reg_OT}, verifies $||\nabla f^* ||=1$ almost surely (\textit{gradient norm preserving} (GNP) architecture~\cite{pmlr-v97-anil19a} ). We apply the approach described in~\cite{pmlr-v97-anil19a}, based on the use of specific activation functions and a process of normalization  of the weights. Two norm preserving activation functions are proposed: i)\ \textbf{GroupSort2} : sorting the vector by pairs, ii) \textbf{FullSort} : sorting the full vector.
These functions are vector-wise rather than element-wise. We also use the P-norm pooling~\cite{boureau2010}, with $P=2$ which is a norm-preserving average pooling.
Concerning linear functions, a weight matrix $W$ is gradient norm preserving if and only if all the singular values of $W$ are equals to $1$. In \cite{pmlr-v97-anil19a}, the authors propose to use the Björck orthonormalization algorithm \cite{bjorck71Ortho}. This algorithm is fully differential and, as for spectral normalization, is applied during the forward inference, and taken into account for back-propagation (see Appendix~\ref{app:normpreserving} for details). We don't consider BCOP \cite{qian_l2-nonexpansive_2019}, which performs slightly better that Björck for convolution but at an higher computation cost. We developed a full tensorflow~\cite{tensorflow2015-whitepaper} implementation in an opensource library, called \Deellip, that enables training of $k$-Lipschitz neural networks, and exporting them as standard layers for inference. 
\subsection{Multi-class hinge-KR classifier}
To adapt our approach to the multi-class case, we propose to use $q$ binary one-vs-all classifiers, where $q$ is the number of classes. The set of labels is now $Y=\{C_1,\ldots,C_q\}$. We name $P_k=\mathbb{P}(X|Y= C_k)$ and $\neg P_k=\mathbb{P}(X|Y\not = C_k)$ the  conditional  distributions  with  respect  to  $Y$. We obtain the following optimization problem :
\begin{equation}
\label{eq:hkr_multi}
\begin{split}
&\underset{f_1,\ldots,f_q \in Lip_1(\Omega)}{\text{sup}}-\mathcal{L}^{hKR}_\lambda(f_1,\ldots,f_q) \\
&\text{where}\\
&\mathcal{L}^{hKR}_\lambda(f_1,\ldots,f_q) = \sum_{k=1}^q \left[\underset{\textbf{x}  \sim \neg P_k}{\mathbb{E}} \left[f_k(\textbf{x})\right] -\underset{\textbf{x} \sim P_k}{\mathbb{E}}\left[f_k(\textbf{x})\right]   +\lambda\underset{\textbf{x}}{\mathbb{E}} \left(m-(2*\underset{y=C_k}{\mathbbm{1}} -1).f_k(x)  \right)_+ \right].
\end{split}
\end{equation}
The global architecture is the same as the binary one except that the last layer has $q$ outputs. For this last layer, each weights corresponding to each output neuron are normalized independently creating $q$ 1-Lipschitz functions with gradient norm preserving properties. With this architecture, the optimal transport interpretation is still valid. The class predicted corresponds to the argmax of the classifier functions. With this approach, the provable robustness lower-bound is the half of the difference between the max and the second max values of $\{f_1(x),\ldots,f_q(x)\}$.  

In \cite{li2019preventing}, Li et al. use classical multi-class hinge loss based on un-centered margin and apply constraint on the entire last layers rather than doing it independently. It allows to a lower provable adversarial robustness bound than ours. However, the $f_1,\ldots,f_q$ obtain in this setting are no more 1-Lipschitz one by one, making the adversarial robustness bounds not comparable directly (see Appendix~\ref{sec:robustness_bounds}).